\chapter{Python の基本}
\label{chap:basics}

\section{データ型と基本操作}

ターミナルで \code{python3} と入力すると、Python をインタラクティブモードで起動できる。
インタラクティブモードに入ると、ターミナルに以下のような文字列が表示される。

\begin{lstlisting}[language=Python, numbers=none, xleftmargin=2\zw]
$ python3
Python 3.13.5 (main, Jun 11 2025, 15:36:57) [Clang 17.0.0 (clang-1700.0.13.3)] on darwin
Type "help", "copyright", "credits" or "license" for more information.
>>>
\end{lstlisting}

この章の前半には、インタラクティブモードへの入力と、その直後に表示される値を対にして示す。

\subsection{変数と基本型}\index{へんすう@変数|textbf}\index{かた@型|textbf}\index{せいすう@整数}\index{ふどうしょうすうてんすう@浮動小数点数}\index{もじれつ@文字列}\index{しんぎち@真偽値}

Python では、値に名前を付けて使う。たとえばデータを表すいくつかの変数を定義すると次のようになる。

\begin{lstlisting}[language=Python, numbers=none, xleftmargin=2\zw]
>>> n_rows = 12
>>> threshold = 3600.0
>>> label = "A"
>>> use_log = True
\end{lstlisting}

\code{n\_rows} は整数、\code{threshold} は浮動小数点数、\code{label} は文字列、\code{use\_log} は真偽値である。実際に \code{type()}\index{type@\code{type()}|textbf} で確かめると次のようになる。解析では、整数と浮動小数点数を区別して扱う感覚が重要になる。たとえば、件数は整数だが、フィットで得られるパラメータは浮動小数点数であることが多い。

\begin{lstlisting}[language=Python, numbers=none, xleftmargin=2\zw]
>>> type(n_rows)
<class 'int'>
>>> type(threshold)
<class 'float'>
>>> type(label)
<class 'str'>
>>> type(use_log)
<class 'bool'>
\end{lstlisting}

Python では、変数名に型を宣言してから値を代入する必要はない。ただし、型がなくなるわけではなく、\code{12}、\code{12.0}、\code{"12"} という三つの値は、それぞれ \code{int}、\code{float}、\code{str} として区別される。このように値が実行時に型を持つ仕組みを動的型付けという。柔軟に書ける一方、大量の数値を一つずつ Python オブジェクトとして処理すると、固定長の数値配列に比べて管理情報や演算呼出しの負担が大きくなる。

後述する NumPy などを用いて大量のデータを高速に処理する際には、C 言語と同様に、メモリ上の表現まで区別した型を指定する場面がある。Python の型と NumPy の型の違いは第~\ref{sec:numpy-types}~節で説明する。基本構文の段階でも、\code{1} は整数、\code{1.0} は浮動小数点数であり、\code{"12"} は数値ではなく文字列であることを区別する。特に文字列の \code{"12"} と整数の \code{12} の混同は、ファイルの読込みや数値計算でエラーの原因になる。

たとえば、見た目が似ていても次の 3 つは別物である。

\begin{lstlisting}[language=Python, numbers=none, xleftmargin=2\zw]
>>> count = 12
>>> threshold = 12.0
>>> text = "12"
>>>
>>> type(count), count
(<class 'int'>, 12)
>>> type(threshold), threshold
(<class 'float'>, 12.0)
>>> type(text), text
(<class 'str'>, '12')
\end{lstlisting}

整数型は件数や添字に、浮動小数点型は平均値やフィット係数などの連続量に用いる。一方、\code{"12"} は文字列であり、そのままでは数値計算に使えない。ファイルから読み込んだ値は文字列になっている場合があるため、見た目が数字でも型を確認する必要がある。

また、\code{=} は「左辺と右辺が等しい」という意味ではなく、「右辺の値を左辺の変数名へ結び付ける」という代入の記号である。条件の比較に用いる \code{==} とは役割が異なる。条件式については、第~\ref{subsec:if}~節で詳しく説明する。

\subsection{数値と演算子}\index{えんざんし@演算子|textbf}\index{しそくえんざん@四則演算}\index{べきじょう@べき乗}\index{じょうよ@剰余}\index{せいすうじょざん@整数除算}

Python は、対話的な計算機として使い始めるのが理解しやすい。ターミナルで \code{python3} と打つと、Python を対話モードで使用できる。
四則演算、べき乗、剰余、整数除算など、基本的な計算の書き方は次の通りである。

\begin{lstlisting}[language=Python, numbers=none, xleftmargin=2\zw]
$ python3
>>> 2 + 3
5
>>> 2 * 3
6
>>> 7 / 3
2.3333333333333335
>>> 7 // 3
2
>>> 7 % 3
1
>>> 2 ** 10
1024
\end{lstlisting}

\code{/} は通常の除算を行い、整数どうしでも浮動小数点数を返す。\code{//} は商を負の無限大方向へ切り下げる床除算であり、正の整数どうしでは小数部分を切り捨てた結果と一致する。観測時間を区間へ分ける、ビン番号を求める、といった処理では、この違いが重要になる。特に負の値を含む場合、\code{-7 // 3} は \code{-3} であり、0 方向へ切り捨てる \code{int(-7 / 3)} の \code{-2} とは一致しない。演算の優先順位は一般的な数式とほぼ同じであり、掛け算や割り算が足し算や引き算より先に評価される。読み手に迷いを残したくないときは、括弧を使って意図を明示した方が安全である。


\begin{lstlisting}[language=Python, numbers=none, xleftmargin=2\zw]
>>> 2 + 3 * 4
14
>>> (2 + 3) * 4
20
>>> 7 / 3
2.3333333333333335
>>> 7 // 3
2
>>> int(7 / 3)
2
\end{lstlisting}

この例では、\code{7 // 3} と \code{int(7 / 3)} はどちらも \code{2} になるが、計算規則は異なる。\code{//} は床除算であり、\code{int(7 / 3)} はまず浮動小数点数の除算を行い、その結果を 0 方向へ切り捨てて整数へ変換している。正の値だけでは差が見えないため、負の値を含む処理では両者を混同してはならない。

浮動小数点数の計算では、表示された値が見た目通りに内部で表現されているとは限らない。これは Python 特有の問題ではなく、一般的な計算機の表現上の性質である。後半でフィットや統計量を扱うときにも、この前提を意識しておきたい。

\subsection{文字列}\index{もじれつ@文字列|textbf}\index{fもじれつ@f文字列}

ファイル名、日付時刻、ヘッダ、ラベルなど、解析コードでは文字列を扱う場面が多い。
Python では、\code{"} や \code{'} で挟まれたものを文字列として認識する。
文字列は単なる表示用のデータではなく、分割、検索、整形の対象でもある。短い操作は、対話モードで 1 つずつ確かめると理解しやすい。

\begin{lstlisting}[language=Python, numbers=none, xleftmargin=2\zw]
>>> name = "detector_A"
>>> date = "2026-04-04"
>>> raw_line = "  detector_A,42,3.5  \n"
>>> name.upper()
'DETECTOR_A'
>>> date.replace("-", "/")
'2026/04/04'
>>> name.split("_")
['detector', 'A']
>>> raw_line.strip()
'detector_A,42,3.5'
>>> f"{name}_{date}.txt"
'detector_A_2026-04-04.txt'
\end{lstlisting}

文字列では、\code{split()}\index{もじれつ@文字列!split@\code{split()}|textbf}、\code{replace()}\index{もじれつ@文字列!replace@\code{replace()}|textbf}、\code{strip()}\index{もじれつ@文字列!strip@\code{strip()}|textbf} のような操作が、ファイル読込みの前処理で頻繁に現れる。特に \code{split()} は、空白区切りやカンマ区切りの 1 行を各列へ分ける最初の道具である。\code{strip()} は行末の改行や前後の空白を落とすときに便利で、テキストファイルを読むとすぐに必要になる。

また、解析コードではファイル名やラベルを組み立てることが多いので、f 文字列\index{もじれつ@文字列!fもじれつ@f文字列|textbf}も早い段階で使えるようにしておくとよい。上の例のように書くと、変数の値を埋め込んだ文字列を簡潔に作れる。ログ出力、図の保存名、結果ファイル名などで頻繁に使うので、単なる飾りではなく実務的な道具として覚えておきたい。

文字列は単なる表示用のデータではなく、検索や分割や整形の対象である。NumPy や pandas を使う前に、この感覚を持っておくとデータ形式の理解が速い。

\section{リスト・タプル・辞書・集合}

\subsection{リスト、タプル、辞書}\index{りすと@リスト|textbf}\index{たぷる@タプル|textbf}\index{じしょ@辞書|textbf}\index{じしょ@辞書!きー@キー}\index{じしょ@辞書!あたい@値}

複数の値をまとめるための基本的なコンテナがある。よく使う 3 種類の定義例を次に示す。

\begin{lstlisting}[language=Python, numbers=none, xleftmargin=2\zw]
>>> labels = ["A", "B", "C", "D"]
>>> point = (12, 18)
>>> row = {"id": 42, "value": 3.5, "flag": True}
\end{lstlisting}

リストは順序付きで変更可能、タプルは順序付きで変更不可、辞書は名前付きの値を持つ。実際に要素を取り出すときは、リストやタプルでは位置を、辞書ではキーを使う。対話モードで試すと次のようになる。

\begin{lstlisting}[language=Python, numbers=none, xleftmargin=2\zw]
>>> labels[0]
'A'
>>> point[1]
18
>>> row["id"]
42
>>> row["value"]
3.5
\end{lstlisting}

この例では、\code{labels[0]} は最初のラベル、\code{point[1]} は 2 番目の座標成分、\code{row["id"]} は \code{"id"} というキーに対応する値を表している。最初のうちは辞書が読みやすいが、大きなデータを大量に処理するときはオーバーヘッドが大きくなることもある。この違いは後の高速化章で再び出てくる。

どの入れ物を選ぶかは、見た目の好みではなく、何をしたいかで決まる。位置でアクセスするならリストやタプル、名前でアクセスするなら辞書が自然である。小さなコードでは差が見えにくいが、処理量が増えると選択の重みが増してくる。

変更できるかどうかの違いも、短い例で見ておくと分かりやすい。リストはあとから要素を書き換えられる。
\begin{lstlisting}[language=Python, numbers=none, xleftmargin=2\zw]
>>> labels[1] = "B_detector"
>>> labels
['A', 'B_detector', 'C', 'D']
\end{lstlisting}

一方、タプルは作った後で要素を書き換えられない。たとえば次のようにするとエラーになる。

\begin{lstlisting}[language=Python, numbers=none, xleftmargin=2\zw]
>>> point[0] = 15
TypeError: 'tuple' object does not support item assignment
\end{lstlisting}

座標のように「ひとかたまりの値として持ち、途中で書き換えない」ものにはタプルが向いている。逆に、ラベル一覧や処理対象のファイル一覧のように、あとから追加・削除・更新したいものにはリストが自然である。辞書は、列名や設定名のように「名前で値を引きたい」場面で特に読みやすい。

\subsection{リストのメソッド}\index{めそっど@メソッド|textbf}\index{append@\code{append()}|textbf}\index{sort@\code{sort()}|textbf}\index{sorted@\code{sorted()}|textbf}\index{pop@\code{pop()}|textbf}

リストにはいくつかの基本メソッドがある。これらを知っていると、簡単な前処理や集約を自然に書ける。主な操作は次のようになる。

\begin{lstlisting}[language=Python, numbers=none, xleftmargin=2\zw]
>>> values = [3, 1, 4]
>>> values
[3, 1, 4]
>>> values.append(1)
>>> values
[3, 1, 4, 1]
>>> values.sort()
>>> values
[1, 1, 3, 4]
>>> last = values.pop()
>>> last
4
>>> values
[1, 1, 3]
\end{lstlisting}

\code{append()} は末尾へ追加し、\code{sort()} はその場で並べ替え、\code{pop()} は末尾要素を取り出して削除する。

リストメソッドの多くは「その場で変更する」点に注意が必要である。新しいリストを返すのではなく、元のリスト自体を書き換える操作が多いので、どこで状態が変わるかを意識して読む必要がある。

特に \code{sort()} は、並べ替えたリストを返すのではなく、元のリストそのものを変更する。そのため、次のように書くと意図とずれることがある。
\begin{lstlisting}[language=Python, numbers=none, xleftmargin=2\zw]
>>> values = [3, 1, 4]
>>> result = values.sort()
>>> result
None
>>> values
[1, 3, 4]
\end{lstlisting}

新しい並べ替え済みリストがほしいときは、\code{sorted()} を使う方が分かりやすい。
\begin{lstlisting}[language=Python, numbers=none, xleftmargin=2\zw]
>>> values = [3, 1, 4]
>>> new_values = sorted(values)
>>> values
[3, 1, 4]
>>> new_values
[1, 3, 4]
\end{lstlisting}

\subsection{インデックスとスライス}\index{いんでっくす@インデックス|textbf}\index{そえじ@添字|see{インデックス}}\index{すらいす@スライス|textbf}\index{ふのそえじ@負の添字}

解析では、配列や文字列の一部を取り出す操作が頻出する。Python では、\code{[]} の中に位置を指定して要素を取り出す。この位置を表す整数をインデックスとよぶ。\code{:} を用いて範囲を切り出す操作はスライスとよばれる。

\begin{lstlisting}[language=Python, numbers=none, xleftmargin=2\zw]
>>> values = [10, 20, 30, 40, 50]
>>>
>>> values[0]
10
>>> values[-1]
50
>>> values[1:4]
[20, 30, 40]
>>> values[:3]
[10, 20, 30]
>>> values[2:]
[30, 40, 50]
>>> values[:-1]
[10, 20, 30, 40]
\end{lstlisting}

リストのインデックスは \code{0} から始まる。負のインデックスは末尾から数え、\code{-1} は最後の要素を表す。
この例では、\code{values[1:4]} はインデックス 1 から 3 までの要素、\code{values[:3]} は \code{values[3]} より前の要素、\code{values[2:]} は \code{values[2]} 以降の要素を取り出す。
リストのスライスでは、\code{:} の左側に指定した位置の要素は含まれるが、右側に指定した位置の要素は含まれない。

文字列に対しても同じ考え方が使える。例えば次のようになる。
\begin{lstlisting}[language=Python, numbers=none, xleftmargin=2\zw]
>>> hhmmss = "123045"
>>> hhmmss[0:2]
'12'
>>> hhmmss[2:4]
'30'
>>> hhmmss[4:6]
'45'
\end{lstlisting}

このように、\code{hhmmss[0:2]} で時、\code{hhmmss[2:4]} で分、\code{hhmmss[4:6]} で秒を取り出せる。この一貫性が Python の扱いやすさの一つである。

\section{条件分岐と反復処理}

\subsection{条件分岐}\index{じょうけんぶんき@条件分岐|textbf}\index{if@\code{if}|textbf}\index{elif@\code{elif}}\index{else@\code{else}}\index{ひかくえんざんし@比較演算子}\index{ろんりえんざんし@論理演算子}
\label{subsec:if}

条件によって処理を変えるには \code{if} を使う。最も基本的な形は次の通りである。

\begin{lstlisting}[language=Python, numbers=none, xleftmargin=2\zw]
>>> value = 2.5
>>> threshold = 3.0
>>> if value <= threshold:
...     print("small")
... else:
...     print("large")
...
small
\end{lstlisting}

比較に使う主な演算子には、\code{==}、\code{!=}、\code{<}、\code{<=}、\code{>}、\code{>=} がある。たとえば \code{value <= threshold} は「\code{value} が \code{threshold} 以下なら真」を意味する。条件式そのものの値は真偽値であり、\code{True} か \code{False} のどちらかになる。
\begin{lstlisting}[language=Python, numbers=none, xleftmargin=2\zw]
>>> value = 5
>>> threshold = 3
>>> value > threshold
True
>>> value == threshold
False
\end{lstlisting}

条件が 3 通り以上に分かれるときは、\code{elif} を使うと読みやすい。たとえば値の大小を 3 段階に分けるなら次のように書ける。

\begin{lstlisting}[language=Python, numbers=none, xleftmargin=2\zw]
>>> value = 3.0
>>> threshold = 3.0
>>> if value < threshold:
...     print("small")
... elif value == threshold:
...     print("equal")
... else:
...     print("large")
...
equal
\end{lstlisting}

解析コードでは、閾値判定、ファイル名の選別、欠損値の除外など、条件分岐が頻繁に現れる。重要なのは、「なぜその条件で分けているのか」を自分で説明できることだ。

条件分岐が増えて読みにくくなったときは、条件式そのものを変数へ切り出すと意図が明確になる。たとえば \code{is\_valid\_row = value <= threshold} のように名前を付けると、後続の条件分岐を短く記述できる。

条件式は、左から右へ読んだときに意味を追える形にするとよい。複数の条件を同時に満たす場合には \code{and}、少なくとも一方を満たす場合には \code{or}、真偽を反転する場合には \code{not} を用いる。

\begin{lstlisting}[language=Python, numbers=none, xleftmargin=2\zw]
>>> value = 5
>>> lower = 3
>>> upper = 10
>>> use_log = False
>>> lower <= value and value <= upper
True
>>> value < lower or value > upper
False
>>> not use_log
True
\end{lstlisting}

否定が重なる条件や、複数の \code{and} と \code{or} が混ざる条件は、括弧や一時変数で整理した方が安全である。解析コードでは、条件の誤読がそのままバグになる。

\subsection{反復処理}

\code{for} 文は、複数の値を順に取り出して同じ処理を適用する構文である。

\begin{lstlisting}[language=Python, numbers=none, xleftmargin=2\zw]
>>> for label in labels:
...     print(label)
...
A
B_detector
C
D
\end{lstlisting}

要素の番号（添字）も必要なら \code{enumerate} を組み合わせると便利である。

\begin{lstlisting}[language=Python, numbers=none, xleftmargin=2\zw]
>>> for i, value in enumerate(labels):
...     print(i, value)
...
0 A
1 B_detector
2 C
3 D
\end{lstlisting}

一方、大規模なデータでは、Python レベルの \code{for} ループが速度の瓶頸になることがある。その回避方法は本書の後半で扱う。

初学者の段階では、まず正しく読める \code{for} 文を書くことが先である。速く書くことを意識するのは、その後でよい。読みやすい基本形を知っているからこそ、後で NumPy や pandas へ置き換えたときに、何を省略しているかが分かる。

\subsubsection{\code{range} と \code{while}}

\code{for} 文と並んで、反復処理では \code{range} と \code{while} もよく用いる。\code{range} は連続した整数を順に生成し、\code{while} は条件が成り立つ間くり返す。記述例を次に示す。

\begin{lstlisting}[language=Python, numbers=none, xleftmargin=2\zw]
>>> for i in range(5):
...     print(i)
...
0
1
2
3
4
>>> count = 0
>>> while count < 3:
...     print(count)
...     count += 1
...
0
1
2
\end{lstlisting}

\code{range(5)} は \code{0} から \code{4} までを返す。ファイル番号やビン番号のような添字処理では便利である。一方、\code{while} は終了条件を自分で管理する必要があるので、無限ループにしないよう注意が必要である。時刻順に読み進めて、条件を満たす間だけ比較するような処理では \code{while} が自然な場合もある。

\subsubsection{\code{break}, \code{continue}, \code{else}}

反復処理では、途中で打ち切ったり、ある条件だけ飛ばしたりしたくなる。Python では \code{break} と \code{continue} がそのために使える。

\begin{lstlisting}[language=Python, numbers=none, xleftmargin=2\zw]
>>> values = [-1, 1, 3, 5]
>>> threshold = 4
>>> for value in values:
...     if value < 0:
...         continue
...     if value > threshold:
...         break
...     print(value)
...
1
3
\end{lstlisting}

\code{continue} はその回だけを飛ばし、\code{break} はループ全体を終了する。探索範囲を打ち切る処理では特に重要である。さらに Python には \code{for ... else} という構文もあり、\code{break} せずに最後まで回り切ったときだけ \code{else} 節が実行される。頻出ではないが、見かけても驚かないようにしておくとよい。

\subsubsection{内包表記}\index{りすとないほうひょうき@リスト内包表記|textbf}

短い変換や簡単なフィルタであれば、内包表記を使うとすっきりと構築できる。

\begin{lstlisting}[language=Python, numbers=none, xleftmargin=2\zw]
>>> values = [1, 2, 3, 4, 5]
>>> squares = [x * x for x in values]
>>> large_values = [x for x in values if x >= 3]
>>> squares
[1, 4, 9, 16, 25]
>>> large_values
[3, 4, 5]
\end{lstlisting}

内包表記は便利だが、何段階もの条件や複雑な変換を 1 行に押し込むと読みにくくなる。初学者のうちは、読みやすさを失わない範囲で使う方がよい。長くなりそうなら、普通の \code{for} 文や補助関数へ戻した方が親切である。

\section{関数とモジュール}

\subsection{関数を分ける理由}\index{かんすう@関数|textbf}\index{かんすう@関数!ひきすう@引数}\index{かんすう@関数!かえりち@返り値}

同じ処理を何度も書かないためだけでなく、処理の意味を分けるためにも関数は重要である。一例として、時刻を表す文字列から秒数を計算する関数を考える。

\begin{lstlisting}[language=Python]
def hhmmss_to_seconds(hhmmss):
    hour = int(hhmmss[0:2])
    minute = int(hhmmss[2:4])
    second = int(hhmmss[4:6])
    return hour * 3600 + minute * 60 + second
\end{lstlisting}

この関数は短いが、「文字列で与えられた時刻を秒へ変換する」という役割が明確である。ここから先は、複数行の処理をまとまりとして読むため、短いスクリプトをファイルへ保存して実行する形も使う。たとえば次の内容を \code{time\_conversion.py} という名前で保存すれば、ターミナルで \code{python3 time\_conversion.py} と実行できる。

\begin{lstlisting}[language=Python, title=time\_conversion.py]
def hhmmss_to_seconds(hhmmss):
    hour = int(hhmmss[0:2])
    minute = int(hhmmss[2:4])
    second = int(hhmmss[4:6])
    return hour * 3600 + minute * 60 + second

print(hhmmss_to_seconds("013015"))
\end{lstlisting}

\begin{lstlisting}[numbers=none, xleftmargin=2\zw]
$ python3 time_conversion.py
4515
\end{lstlisting}

このように、関数は入力を受け取り、決まった規則で出力を返す小さな部品として読むと分かりやすい。解析コードでは、このような小さな関数が積み重なって全体の読みやすさを作る。

関数を切り出す利点は再利用だけではない。入力が何で、出力が何かをはっきりさせることで、テストしやすくなる。短い補助関数ほど後回しにされがちだが、解析の正しさはこうした小さな部品の積み重ねで決まることが多い。

\subsubsection{関数の引数}\index{かんすう@関数!ひきすう@引数|textbf}\index{でふぉるとひきすう@デフォルト引数|textbf}\index{きーわーどひきすう@キーワード引数|textbf}

関数を書くときは、引数の意味をはっきりさせることが重要である。Python では位置引数だけでなく、デフォルト値付き引数やキーワード引数も使える。

\begin{lstlisting}[language=Python]
def histogram_path(name, suffix=".pdf"):
    return f"{name}{suffix}"

print(histogram_path("value_hist"))
print(histogram_path("value_hist", suffix=".png"))
\end{lstlisting}
\begin{lstlisting}[numbers=none, xleftmargin=2\zw]
value_hist.pdf
value_hist.png
\end{lstlisting}

デフォルト引数を使うと、よく使う設定を省略できる一方で、必要なら上書きできる。コマンドライン引数と同じで、解析コードの柔軟性は「固定値をどこに置くか」で決まることが多い。関数でも同じ発想が使える。

ただし、可変オブジェクトをデフォルト引数に置くと予想外の共有が起きる。この罠は Python 公式チュートリアル~\cite{python_tutorial} でも強調されている典型例である。初学者の段階では、リストや辞書をデフォルトにしたくなったら一度立ち止まった方がよい。以下の書き方は避けるべきである。

\begin{lstlisting}[language=Python]
def bad_append(value, buffer=[]):
    buffer.append(value)
    return buffer
\end{lstlisting}

この例では、\code{buffer} が呼び出しのたびに新しく作られるわけではない。そのため、前回の状態が次回にも残る。実際に呼ぶと次のようになる。

\begin{lstlisting}[language=Python]
print(bad_append(1))
print(bad_append(2))
\end{lstlisting}

\begin{lstlisting}[numbers=none, xleftmargin=2\zw]
[1]
[1, 2]
\end{lstlisting}

安全に書くなら \code{None}\index{none@\code{None}|textbf} を使って関数内部で初期化する方がよい。

\begin{lstlisting}[language=Python]
def good_append(value, buffer=None):
    if buffer is None:
        buffer = []
    buffer.append(value)
    return buffer
\end{lstlisting}

\begin{lstlisting}[language=Python]
print(good_append(1))
print(good_append(2))
\end{lstlisting}

\begin{lstlisting}[numbers=none, xleftmargin=2\zw]
[1]
[2]
\end{lstlisting}

初学者がつまずきやすいので、早めに知っておく価値が高い。

\subsubsection{モジュールと \code{import}}\index{もじゅーる@モジュール|textbf}\index{import@\code{import}|textbf}

別ファイルに書かれた関数を使うには次のように \code{import} する。たとえば標準ライブラリの \code{math} を使うと、円周率や平方根を次のように扱える。

\begin{lstlisting}[language=Python]
import math
import glob
import argparse

print(math.pi)
print(math.sqrt(2.0))
\end{lstlisting}

\begin{lstlisting}[numbers=none, xleftmargin=2\zw]
3.141592653589793
1.4142135623730951
\end{lstlisting}

Python の解析では、標準ライブラリと外部ライブラリを組み合わせて使う。各名前の定義元を確認すると、規模の大きいコードでも関数やクラスの役割を追いやすい。

特に、同じ名前に見えても標準ライブラリと外部ライブラリで役割が大きく異なることがある。インポート文は単なるおまじないではなく、そのファイルが何に依存しているかを示す地図でもある。ここまでは既存のライブラリを読み込む話だったが、同じ仕組みで自分が書いたファイルも読み込める。

\subsubsection{自作モジュールへの分離}

解析が少し大きくなると、1 つのファイルへ全部を書くのはすぐに苦しくなる。よく使う関数や定数は別ファイルへ分け、モジュールとして読み込む方が見通しが良い。たとえば次のようなファイル構造で分けるとする。

\begin{lstlisting}[numbers=none, xleftmargin=2\zw]
analysis_project/
|-- constants.py
|-- physics.py
`-- main.py
\end{lstlisting}

この構成にすると、\code{physics.py} に書いた関数を \code{main.py} から読み込める。次の例は、その最小形である。

\begin{lstlisting}[language=Python]
# physics.py
import numpy as np

def deg_to_rad(degrees):
    return degrees * np.pi / 180.0


# main.py
import physics

theta = physics.deg_to_rad(45.0)
\end{lstlisting}

この形にしておくと、座標変換、時刻変換、定数定義のような再利用しやすい処理を 1 か所へまとめられる。たとえば \code{main.py} 側で \code{print(theta)} と書けば、結果は次のようになる。

\begin{lstlisting}[numbers=none, xleftmargin=2\zw]
0.7853981633974483
\end{lstlisting}

ファイルを分けると、「他のファイルから読み込んで使う」場面が自然に増える。逆に、そのファイル自身を直接実行して動作確認したいこともある。次に示す \code{if \_\_name\_\_ == "\_\_main\_\_":} は、その 2 つを切り分けるための書き方である。

\subsubsection{スクリプトとしての実行}

Python ファイルは、他のファイルから読み込まれることもあれば、直接実行されることもある。その違いを区別する基本形が次である。

\begin{lstlisting}[language=Python]
def main():
    print("run")

if __name__ == "__main__":
    main()
\end{lstlisting}

この形にしておくと、モジュールとして読み込んだときには関数だけを使え、直接実行したときだけ本体が動く。後半の解析スクリプトでも、この形を基本として採用すると整理しやすい。

\section{例外・docstring・命名・設定値}

\subsection{例外名とトレースバックによる原因の特定}\index{れいがい@例外|textbf}\index{とれーすばっく@トレースバック|textbf}\index{filenotfounderror@\code{FileNotFoundError}}\index{valueerror@\code{ValueError}}\index{typeerror@\code{TypeError}}

初心者のうちは、エラーメッセージが出るとそこで止まってしまいがちである。しかし、トレースバックは「どこで何が起きたか」を教えてくれる重要な情報である。

たとえば \code{FileNotFoundError} ならファイルパスが違うかもしれないし、\code{ValueError} なら想定した形式でデータを読めていないかもしれない。実際のエラー表示は次のようになる。

\begin{lstlisting}[language=Python]
with open("missing.txt") as f:
    text = f.read()
\end{lstlisting}

\begin{lstlisting}[numbers=none, xleftmargin=2\zw]
Traceback (most recent call last):
  File "read_missing_file.py", line 1, in <module>
    with open("missing.txt") as f:
FileNotFoundError: [Errno 2] No such file or directory: 'missing.txt'
\end{lstlisting}

大切なのは、英語の文章全体を恐れるのではなく、例外名と最後の数行を見ることである。

解析コードでは、開発中にエラーが発生すること自体は珍しくない。問題は、エラーを無視したり、意味を確かめずに場当たり的な修正を加えたりすることである。例外名、発生行、直前に読み込んだデータの形を順に確認すると、原因を切り分けやすい。

\subsubsection{docstring とヘルプ}\index{docstring@docstring|textbf}\index{help@\code{help()}|textbf}

関数やモジュールの役割を後から思い出せるようにするには、短い説明を書いておくとよい。Python では次のように、関数定義の直後に文字列を書くと \code{docstring}（ドキュメント文字列）になる。

\begin{lstlisting}[language=Python]
def read_table(path):
    """空白区切りの表を読み、各行を辞書として返す。"""
    ...
\end{lstlisting}

\code{help(read\_table)} を実行すると、この説明を後から確認できる。たとえば要約だけの docstring なら、次のように表示される。

\begin{lstlisting}[language=Python]
help(read_table)
\end{lstlisting}

\begin{lstlisting}[numbers=none, xleftmargin=2\zw]
Help on function read_table:

read_table(path)
    空白区切りの表を読み、各行を辞書として返す。
\end{lstlisting}

コメントと同じように見えるが、実行時にも参照できる点が異なる。解析コードが増えてくると、どの関数が何を返すのかを書いておくことの価値が大きくなる。

短い関数なら 1 行の docstring で十分だが、少し長い処理では複数行にして、少なくとも「何をするか」「主要な引数」「返り値」「起こり得る例外」を書いておくと読みやすい。たとえば次のような形である。

\begin{lstlisting}[language=Python]
def read_table(path):
    """空白区切りの表を読み込む。

    Args:
        path: 入力ファイルへのパス。

    Returns:
        各行を辞書として格納したリスト。

    Raises:
        FileNotFoundError: ファイルが存在しない場合。
    """
    ...
\end{lstlisting}

流儀はいくつかあるが、最初の 1 行を「要約」、その後ろを詳細と考えると整理しやすい。特に、他人が \code{help()} やエディタ上の補完で関数を読むときは、最初の 1 行がそのまま関数の説明になることが多い。

\subsubsection{名前付けと PEP 8}\index{pep 8@PEP 8|textbf}\index{snake_case@\code{snake\_case}}\index{camelcase@\code{CamelCase}}\index{upper_case@\code{UPPER\_CASE}}

Python には PEP 8 という代表的なスタイルガイドがある。変数名、関数名、空白の入れ方をそろえるだけでも、コードの構造を追いやすくなる。

\begin{lstlisting}[language=Python]
MAX_EVENTS = 1000

def read_event_table(path):
    total_count = 0
    ...
\end{lstlisting}

一般には、変数名や関数名は \code{snake\_case}、定数は大文字の \code{UPPER\_CASE} にすることが多い。さらに、演算子の前後に空白を入れる、1 行を詰め込みすぎない、といった基本だけでも効果は大きい。スタイルの統一は飾りではなく、バグを見つけやすくするための道具だと考える方が実践的である。

PEP 8 を最初から全部暗記する必要はないが、少なくとも次の規則は守る価値が高い。

\begin{itemize}
\item タブではなく 4 スペースによるインデント
\item 関数名・変数名には \code{snake\_case}、クラス名には \code{CamelCase}、定数には \code{UPPER\_CASE} という命名規則
\item 演算子（\code{=}, \code{+}, \code{-} など）の前後と、カンマの後ろの空白
\item トップレベルの関数定義・クラス定義を区切る空行
\item 過度に長い行の回避（PEP 8 の目安は 79 文字）
\item 標準ライブラリ、外部ライブラリ、自作モジュールの順にまとめた \code{import}
\end{itemize}

これらの規則だけでも、表記のばらつきを減らせる。すべての規則を機械的に適用することより、同じプロジェクト内で一貫した書式を保つことが重要である。

\subsubsection{マジックナンバーの排除}\index{まじっくなんばー@マジックナンバー|textbf}
\label{sec:magic-number}

コード中に \code{250} や \code{3600} のような数値が突然現れると、その値が何を表すのかを読む側が推測しなければならない。このような、説明なしに裸で現れる数値をマジックナンバーと呼ぶ。これを避けるには、次のように意味のある名前を持たせる。

\begin{lstlisting}[language=Python]
NS_PER_SECOND = 1_000_000_000
COINCIDENCE_WINDOW_NS = 250

dt_ns = second_diff * NS_PER_SECOND + ns_diff
is_coincident = abs(dt_ns) <= COINCIDENCE_WINDOW_NS
\end{lstlisting}

\code{250} をそのまま書くより、\code{COINCIDENCE\_WINDOW\_NS} と書いた方が意味が明確である。\code{3600} なら \code{SECONDS\_PER\_HOUR}、\code{86400} なら \code{SECONDS\_PER\_DAY} といった具合に、意味を名前へ出した方が読みやすい。

ただし、どんな値でもコード内定数にすればよいわけではない。物理定数や単位換算のように定義として固定したい値はコード内の定数に向いている。一方で、入力パス、出力先、閾値、利用する検出器番号のように実行条件として変わり得る値は、コマンドライン引数や設定ファイルへ出した方がよい。値の置き場所そのものが設計である、という意識を持つことが重要である。

\subsubsection{例外処理}\index{れいがい@例外!れいがいしょり@例外処理|textbf}\index{try@\code{try}|textbf}\index{except@\code{except}|textbf}

エラーは読むべき情報だが、予想される失敗を自分で受け止めたい場面もある。例えば、1 行だけ形式が壊れているデータを見つけて記録し、解析全体は続けたいことがある。そのときに使うのが次のような \code{try} と \code{except} の組み合わせである。

\begin{lstlisting}[language=Python]
try:
    value = float(text)
except ValueError:
    print("数値へ変換できませんでした")
\end{lstlisting}

たとえば \code{text = "abc"} のような文字列に対して \code{float(text)} を呼ぶと、\code{ValueError} が起きる。したがって、何でも \code{except:} で受けてしまうと、本来止まるべきバグまで隠してしまう。どの例外を受けるのかを明示し、何を記録して、何を諦めて、何を継続するのかを自分で決めることが重要である。

\section{イベントデータを用いた基本構文の演習}

この章の課題は、次の小さなイベント列だけで実行できる。前章までのファイルは必要ない。

\begin{lstlisting}[language=Python]
events = [
    {"detector": "A", "time_ns": 1000, "signal": 12.4},
    {"detector": "B", "time_ns": 1180, "signal": 9.8},
    {"detector": "A", "time_ns": 1300, "signal": 15.1},
    {"detector": "B", "time_ns": 1510, "signal": None},
]
\end{lstlisting}

\begin{enumerate}
\item \textbf{条件に合う入力行の表示}。
\code{for} 文で各辞書を取り出し、\code{signal} が \code{None} でない行だけを \code{A 1000 12.4} の形式で表示せよ。\code{is not None} を使う理由と、辞書の角括弧に書いたキーの役割を説明せよ。表示される行は 3 行である。

\item \textbf{関数への分離と集計}。
イベント列と検出器名を受け取り、該当する検出器の有効な \code{signal} の平均を返す \code{mean\_signal(events, detector)} を作れ。A では \code{13.75}、B では \code{9.8} が返ることを \code{print()} で確認せよ。該当する有効値がない場合は \code{ValueError} を送出し、docstring に引数、返り値、例外条件を明記せよ。

\item \textbf{時刻差の計算}。
\code{time\_ns} だけをリスト内包表記で取り出し、昇順に並べた後、隣り合う要素の差を求めよ。結果は \code{[180, 120, 210]} になる。\code{range(len(times) - 1)} の開始値、終了値、最後の添字がどう対応するかを説明せよ。

\item \textbf{コマンドライン引数による条件指定}。
検出器名を位置引数、信号の下限を \code{--min-signal} で受け取るスクリプトを \code{argparse} で作れ。既定値は \code{0.0} とする。\code{python filter\_events.py A --min-signal 13} を実行したとき、時刻 \code{1300} の 1 件だけが表示されることを確認し、位置引数とオプション引数の違いを説明せよ。
\end{enumerate}

この章で扱った値、コンテナ、条件分岐、反復、関数は、ファイルから読み取ったデータを処理するための部品になる。次章では、これらの部品を入出力と結び付け、テキスト、CSV、JSON、バイナリ形式から Python オブジェクトへの変換を扱う。
